% ---------------------------------------------------------------------------
% SUPERSEDED. Not the current paper.
% Early Interspeech-style draft, kept for the record. Several of its numbers
% were later corrected (padding fix, measured noise floor, saturation ~240 ms
% -> ~340 ms). The current paper is taslp/taslp.tex, submitted to TASLP.
% ---------------------------------------------------------------------------
% ---------------------------------------------------------------------------
% ICASSP 2027 submission. Deadline 16 Sep 2026, Toronto.
%
% Uses the OFFICIAL kit downloaded from
%   https://cmsworkshops.com/ICASSP2027/papers/PaperFormat/
% namely spconf.sty and IEEEbib.bst, both vendored alongside this file.
%
% HARD LIMIT: 5 pages total. Pages 1-4 carry technical content, figures and
% references; page 5 may contain ONLY references, funding acknowledgements and
% the Compliance with Ethical Standards statement. Reviewing is NOT blind, so
% the author block stays.
%
% The Interspeech/SSW version of this paper is ../main.tex, which is 4+2 and
% therefore fits SSW14 (4 pages + up to 2 for references) without trimming.
% Keep the two in sync by hand; the body below is a copy, not an include,
% because the two venues need different front matter and different trims.
% ---------------------------------------------------------------------------
\documentclass{article}
\usepackage{spconf,amsmath,amssymb,graphicx,booktabs,url}
\usepackage[T1]{fontenc}
\usepackage{tipa}

\newcommand{\blfootnote}[1]{%
  \begingroup\renewcommand\thefootnote{}\footnote{#1}\addtocounter{footnote}{-1}\endgroup}
\newcommand{\talg}{\ensuremath{t_{\mathrm{alg}}}}
\newcommand{\tcmp}{\ensuremath{t_{\mathrm{cmp}}}}
\newcommand{\tbuf}{\ensuremath{t_{\mathrm{buf}}}}

\title{How Much Future Context Does Streaming Accent Conversion Need?\\
A Controlled Study Using Causal Phone Translation}

\name{Dipesh Tripathi}
\address{University of South Dakota \\ \texttt{dipesh.tripathi@coyotes.usd.edu}}

\begin{document}
\ninept
\maketitle


\begin{abstract}
\noindent
Streaming foreign accent conversion (FAC) systems choose a lookahead budget and
report a single end-to-end latency, but no published work sweeps that budget or
separates the delay a modelling choice imposes from the delay the hardware
imposes. We train a causal phone translator at 16 lookaheads from 0 to 640\,ms
with every other factor held fixed, and report four results. First, quality
improves smoothly with lookahead at $-0.045$ phone error rate (PER) per
doubling ($R^2=0.98$) with \emph{no locatable knee}: a piecewise fit is
preferred under every axis treatment we tried, yet the breakpoint moves between
40 and 180\,ms and even the narrowest bootstrap interval spans 2.4 octaves.
Instead of a knee we report a \emph{saturation point} at ${\approx}240$\,ms,
beyond which each additional 20\,ms frame buys less than the measured
$2\sigma$ noise floor. A 40\,ms budget, the value used by the current
state of the art, forgoes ${\approx}40\%$ of the achievable PER reduction.
Second, a pre-registered hypothesis that vowels and approximants degrade faster
than obstruents is \emph{not supported}: all seven manner classes improve by
52--72\% relative, and the pre-registered grouping explains less variance than
the spread within its own groups. It still fails when errors are reweighted by
measured acoustic cost, so this is not an artefact of PER charging every error
equally. Third, canonical-phone conversion benefits more from future context than
accent-faithful transcription at matched capacity ($-0.045$ against $-0.029$
PER per doubling, $1.55\times$), and within the canonical-phone arm the benefit is
$1.32\times$ larger at positions where the speaker actually deviated from
canonical. Fourth, per-chunk compute is flat in
lookahead ($-0.3\%$ to $+1.9\%$ across a $17\times$ change in algorithmic
latency, replicated on two architectures), which makes the two latency terms
independently actionable. We release the harness and the per-condition outputs.
\end{abstract}

\blfootnote{Harness, causality proof, sweep configurations, analysis scripts
with self-tests and per-condition hypothesis dumps:
\url{https://github.com/Dipeshtripathi13/streaming-fac-lookahead}. Data is
L2-ARCTIC~\cite{l2arctic2018} via \texttt{KoelLabs/L2Arctic} (CC-BY-NC-4.0);
that licence is non-commercial, so configurations are released rather than
weights.}

\begin{keywords}
accent conversion, streaming speech, latency, lookahead, self-supervised models
\end{keywords}

\section{Introduction}

Streaming FAC exists and works. PHONOS~\cite{phonos2026} runs at
${\leq}40$\,ms lookahead and reports an 81\% reduction in non-native accent
confidence. What does not exist is any account of \emph{why} 40\,ms, what it
costs in quality, or what any of it does on hardware that is not a datacentre
GPU.

Three gaps follow, and this paper addresses the first two directly.

\textbf{The budget is chosen, not characterised.} PHONOS \emph{states} 40\,ms
and publishes no ablation showing 40\,ms is necessary or sufficient; TVTSyn
fixes 4 future tokens ($\approx$80\,ms) without varying them. DarkStream does
ablate, over three points $\{0, 140, 280\}$\,ms, and finds the third buys
$<$1pp, but three points scored on anonymisation accuracy cannot locate
where returns stop, which is the question a deployment budget actually asks.
``What would 160\,ms buy you'' still has no published answer.

\textbf{Algorithmic and computational latency are reported fused.} A figure
such as ``${<}241$\,ms on a single GPU'' combines a modelling choice with
unstated inference and buffering costs. A reader cannot tell whether a faster
chip would help. We decompose throughout as
\begin{equation}
t_{\mathrm{e2e}} = \underbrace{c + L}_{t_{\mathrm{alg}}} + t_{\mathrm{cmp}} + t_{\mathrm{buf}},
\label{eq:decomp}
\end{equation}
where $c$ is chunk size, $L$ is lookahead, \tcmp{} is per-chunk compute and
\tbuf{} is I/O and jitter buffering. \talg{} is inherent: it is the delay on an
infinitely fast machine, because the model refuses to emit frame $n$ until it
has seen frame $n{+}L$. \tcmp{} is contingent, quantise, prune, buy a faster
chip.

\textbf{Degradation is reported in aggregate.} Existing results are
utterance-level MOS, WER or accentedness. Which \emph{sounds} break first is
unreported.

Our contributions:
\begin{enumerate}\itemsep1pt
\item A latency--quality characterisation of streaming FAC across 16 lookahead
      budgets, reporting an exchange rate and a floor-relative saturation
      point rather than an unidentifiable knee (\S\ref{sec:rq1}).
\item A negative result on a pre-registered phoneme-class hypothesis
      (\S\ref{sec:rq2}).
\item Evidence that accent conversion needs more lookahead than
      accent-faithful transcription at matched capacity (\S\ref{sec:rq3}).
\item Two patches that make a masked SSL encoder actually causal, with a
      truncation proof (\S\ref{sec:causal}).
\end{enumerate}

\section{Related work}

Streaming FAC is established. Nguyen et al.~\cite{nguyen2025streaming} present
a streaming non-autoregressive accent-conversion system on an Emformer encoder,
demonstrating feasibility at stable latency; PHONOS~\cite{phonos2026}
subsequently targets much lower latency with bounded future context. Neither
characterises how quality varies with that context, which is the axis we
isolate. PHONOS, TVTSyn~\cite{tvtsyn2026} and DarkStream~\cite{darkstream2025}
come from one group and pair conversion with speaker anonymisation, and two
report more than is sometimes credited. TVTSyn gives CPU as well as GPU latency
($\approx$132 and $\approx$79\,ms at 60\,ms chunk), so a CPU number is not
unpublished, a CPU latency--\emph{quality} curve is. DarkStream ablates
look-ahead over $\{0,140,280\}$\,ms and selects 140\,ms because the extra
140\,ms buys $<$1pp at double the buffering delay: a three-point ablation,
scored on anonymisation accuracy rather than conversion quality, whose marginal
return collapses in the same region where we place saturation on phone error
rate over 16 budgets. That agreement across objectives is a stronger position
than absent prior evidence would have been.

StreamVC~\cite{streamvc2024} reports mobile operation without device, core or
thread details, limiting hardware-level reproducibility. A
survey~\cite{acsurvey2026} frames the space; streaming
anonymisation~\cite{streamvoiceanon2026, streamvoiceanonplus2026, e2eanon2024}
shares the constraint but not the accent objective; offline
FAC~\cite{zhao2021fac, accentron2022} and low-latency VC~\cite{llvc2023}
bracket the design space. Our contribution is orthogonal: we characterise an
axis these systems parameterise but do not study, and claim no quality
improvement over any of them.

\section{Method}

\subsection{A causal phone translator as the probe}
\label{sec:method-translator}

Training a full FAC pipeline per lookahead is prohibitive. We instead isolate
the component the research question is about. L2-ARCTIC~\cite{l2arctic2018}
provides, per utterance, both the canonical phone sequence a native speaker
would produce (\texttt{g2p}) and the sequence the non-native speaker actually
produced (\texttt{ipa}). Accent conversion is then directly trainable as
\begin{center}
\emph{non-native audio} $\rightarrow$ \emph{native phone sequence},
\end{center}
a CTC head over a frozen causal encoder. This is the supervised core of the
task, it costs a GPU-day rather than 200 GPU-hours, and swapping the target
tensor from \texttt{g2p} to \texttt{ipa} yields a matched-capacity control that
differs in \emph{one tensor}, a tighter control than AC-versus-VC.

\subsection{Making the encoder causal}
\label{sec:causal}

Masking self-attention in wav2vec2/WavLM-base~\cite{wavlm2022, wav2vec2} does
\emph{not} produce a causal encoder. Two leaks remain: the positional
convolution (\texttt{pos\_conv\_embed}, depthwise, kernel 128, symmetric
padding) admits 1.28\,s of future at a 20\,ms frame rate; and the
feature-encoder GroupNorm normalises each channel over the entire utterance,
making every output frame depend on every input frame.

We verify causality by truncation: delete audio after frame $t$ and check
whether any earlier frame moved (Table~\ref{tab:causal}). Note the third row:
patching GroupNorm alone makes matters \emph{worse}. A partial causality fix is
not a partial improvement.

\begin{table}[t]
\centering\small
\caption{Truncation test on \texttt{wavlm-base-plus}. Relative $L_2$ change in
earlier frames after deleting later audio.}
\label{tab:causal}
\begin{tabular}{lrc}
\toprule
configuration & relative $L_2$ & causal? \\
\midrule
attention mask only        & $1.14\times10^{-2}$ & no \\
\quad + causal positional conv & $6.00\times10^{-3}$ & no \\
\quad + cumulative GroupNorm   & $2.60\times10^{-2}$ & no \\
\textbf{both patches}      & $\mathbf{6.14\times10^{-6}}$ & \textbf{yes} \\
\bottomrule
\end{tabular}
\end{table}

\subsection{Lookahead is quantised}

Lookahead is realised as $\lceil L/f \rceil$ frames at frame rate $f=20$\,ms.
A knee narrower than 20\,ms is not merely unmeasured but
\emph{unrepresentable}, and sampling finer than $f$ silently creates duplicate
conditions that reweight any fit. Our dense grid is therefore 20\,ms, and we
assert zero collisions in $t_{\mathrm{alg}}$ across all conditions before
fitting.

\section{Experimental setup}

Frozen causal WavLM-base-plus, CTC phone head, 1200 steps, batch 8, chunk
40\,ms, 100-frame lookback, speaker-disjoint test split of 400 utterances.
Two sweeps: a 7-lookahead $\times$ 2-target $\times$ 3-seed sweep (42 runs)
and a 16-lookahead $\times$ 2-seed dense sweep, run on \emph{both} target arms
(64 runs), all on an NVIDIA T4. Compute measurements use an Apple M4 and an ARM64
Linux host, CUDA-event or \texttt{perf\_counter\_ns} timed, 28--50 repetitions
with three warm-up iterations discarded.

\noindent\textbf{Noise floor.} Two sweep processes briefly ran concurrently by
operator error, repeating six conditions at \emph{identical seed and
configuration}. Those repeats are the only direct measurement of fixed-seed
run-to-run noise available, and a multi-seed design cannot produce it: pairing
within seed cancels seed effects but not this residual. Differences ranged
$1.5\times10^{-4}$ to $6.5\times10^{-3}$ PER, giving $\sigma=0.0022$ and a
$2\sigma$ resolvable threshold of $\mathbf{0.0044}$ PER. Every effect below is
gated on it.

\section{Results}

\subsection{Compute is flat in lookahead}
\label{sec:compute}

Across $L\in[0,640]$\,ms at chunk 40\,ms on the M4, \talg{} spans 40--680\,ms
($17\times$) while median per-chunk compute moves from 43.23\,ms to at most
44.06\,ms, a range of $-0.3\%$ to $+1.9\%$. On ARM64 the span is $0$ to
$+7.6\%$. The two terms of Eq.~\ref{eq:decomp} are therefore independently
actionable, which is the practical justification for reporting them
separately.

\textbf{Independently actionable is not the same as free.} At this
configuration $\tcmp{}\approx43$--44\,ms against a 40\,ms chunk, i.e.\
$\mathrm{RTF}=\tcmp/c\approx1.08>1$: the base encoder does not sustain real
time at 40\,ms chunks on this machine, whatever $L$ is. Eq.~\ref{eq:decomp}
describes the delay of a single chunk and must be read alongside the
steady-state throughput constraint
\begin{equation}
\mathrm{RTF} = \tcmp / c < 1,
\label{eq:rtf}
\end{equation}
which lookahead leaves untouched but chunk size and hardware do not. Changing
$L$ is close to free in compute; being able to run at all is not.

\subsection{RQ1: an exchange rate and a saturation point, not a knee}
\label{sec:rq1}

\begin{figure}[t]
\centering
\includegraphics[width=\columnwidth]{fig_dense_no_knee.png}
\caption{Dense sweep, canonical-phone arm, 16 lookaheads $\times$ 2 seeds.
\textbf{Left:} smooth curvature; the shaded band spans the breakpoint estimates
produced by three defensible treatments of $L{=}0$, none identifiable.
\textbf{Right:} marginal return per added 20\,ms frame against the measured
$2\sigma$ floor; red bars are below it.}
\label{fig:dense}
\end{figure}

Canonical-phone-arm PER falls monotonically from $0.4256$ at $L{=}0$ to $0.1561$ at
$L{=}640$\,ms. Fitting $L\geq20$\,ms on $\log_2 L$ gives
\begin{equation}
\Delta\mathrm{PER} = -0.0452 \text{ per doubling},\quad R^2 = 0.983,
\end{equation}
over 15 points.

\textbf{There is curvature, and it is not localisable.} BIC prefers a
two-segment fit under all three axis treatments we tried
($\Delta\mathrm{BIC}=-22.2$, $-31.2$, $-36.8$), which by the usual
thresholds~\cite{kassraftery1995} is strong evidence. But the breakpoint lands
at 40, 180 and 180\,ms respectively, and every bootstrap 90\% interval spans
2.4 to 3.0 octaves. Sixteen points did not rescue a knee. The reason is
mechanical and worth stating: on a $\log_2(L{+}1)$ axis the $L{=}0
\rightarrow 20$\,ms gap is 4.39 units while every other gap is ${\approx}1.0$,
so a piecewise fit places a breakpoint just past that gap whatever the data
does.

\textbf{What we report instead} is the region beyond which marginal return
falls below the measured floor. The grid is 20\,ms up to 200\,ms and coarser
above it, so past that point we report the \emph{average} PER reduction per
20\,ms equivalent of added lookahead over each sampled interval, not per
individually observed 20\,ms step. On that basis gains run $+0.056$
(0$\rightarrow$20\,ms), $+0.021$ (60$\rightarrow$80), $+0.0045$
(180$\rightarrow$200), and fall below the $0.0044$ floor for every sampled
interval beyond 240\,ms. We therefore report a \textbf{saturation region of
200--280\,ms} rather than a point: the 140$\rightarrow$160 and
180$\rightarrow$200\,ms steps are already only 1.0--1.5$\times$ the floor, the
floor itself is estimated from six repeats, and the intervals above 320\,ms
average over 160\,ms of lookahead. Unlike a knee, this quantity depends on the
floor rather than on an axis choice.

\begin{table}[t]
\centering\small
\caption{Canonical-phone-arm PER against lookahead (mean of 2 seeds) and marginal
return per added 20\,ms frame, in units of the measured $2\sigma{=}0.0044$
floor. Below $1.0\times$ an additional frame is not measurably useful.}
\label{tab:dense}
\begin{tabular}{rrrr}
\toprule
$L$ (ms) & PER & gain/frame & $\times$ floor \\
\midrule
0   & 0.4256 & --     & -- \\
20  & 0.3701 & 0.0556  & 12.7 \\
40  & 0.3275 & 0.0426  & 9.8 \\
60  & 0.2963 & 0.0312  & 7.2 \\
80  & 0.2754 & 0.0209  & 4.8 \\
100 & 0.2564 & 0.0190  & 4.4 \\
120 & 0.2446 & 0.0118  & 2.7 \\
140 & 0.2269 & 0.0176  & 4.0 \\
160 & 0.2206 & 0.0064  & 1.5 \\
180 & 0.2106 & 0.0099  & 2.3 \\
200 & 0.2062 & 0.0045  & 1.0 \\
240 & 0.1973 & 0.0044  & 1.0 \\
280 & 0.1891 & 0.0041  & \textbf{0.9} \\
320 & 0.1823 & 0.0034  & \textbf{0.8} \\
480 & 0.1649 & 0.0022  & \textbf{0.5} \\
640 & 0.1561 & 0.0011  & \textbf{0.3} \\
\bottomrule
\end{tabular}
\end{table}

\textbf{The cost of a 40\,ms budget.} PER at $L{=}40$\,ms is $0.3275$; at
240\,ms it is $0.1973$. \textbf{Relative to the 40\,ms condition, increasing
lookahead to 240\,ms reduces PER by a further 39.8\%}, and to 640\,ms by
52.3\%. Measured instead against the full range available here
($L{=}0\rightarrow240$\,ms), a 40\,ms budget has already captured 57\% of the
attainable reduction and leaves 43\% unclaimed.
``Current budgets are under-provisioned'' is thus supportable as
diminishing-returns geometry with a measurable saturation point, not as a
knee.

\subsection{RQ2: manner class is the wrong axis}
\label{sec:rq2}

We pre-registered, before any of this data existed, that vowels and
approximants would break first (formant trajectories over 80--250\,ms) while
stops, fricatives, affricates and nasals would survive short lookahead (locally
determined gestures). Testing this on 42{,}255 unique reference-phone tokens evaluated across seven
lookahead conditions (295{,}785 phone--condition observations), with each
substitution or deletion charged to the class of the \emph{reference} phone:

\begin{table}[t]
\centering\small
\caption{Relative error-rate reduction, $L{=}0\rightarrow640$\,ms. B =
pre-registered as breaking first, S = as surviving. The predicted grouping does
not order the data.}
\label{tab:h2}
\begin{tabular}{llrr}
\toprule
& class & rel.\ gain & ref.\ phones \\
\midrule
S & nasal        & 0.716 & 4\,020 \\
B & approximant  & 0.700 & 4\,575 \\
S & fricative    & 0.687 & 7\,659 \\
B & vowel (mono) & 0.673 & 13\,008 \\
B & vowel (diph) & 0.654 & 4\,824 \\
S & stop         & 0.615 & 7\,071 \\
S & affricate    & 0.525 & 1\,098 \\
\bottomrule
\end{tabular}
\end{table}

Group means differ in the predicted direction ($0.676$ vs $0.636$), but that is
the weakest available reading and three stronger checks fail. The effect size
is $0.66$: the between-group difference ($0.040$) is smaller than the pooled SD
across classes ($0.060$). The surviving group's \emph{internal} spread is
$0.190$, $4.7\times$ the between-group difference. And 5 of 12 pairwise
orderings are violated, with nasals out-gaining all three breaks-first classes.
\textbf{H2 is not supported.} The finding is not that lookahead fails to help
vowels, every class improves 52--72\% relative, all above the floor. The
finding is that manner class does not predict which sounds need right context.

\textbf{Reweighting each error by its acoustic cost does not rescue it.} PER
charges every error 1.0, so the natural objection is that a substituted vowel
and a substituted stop cannot sound equally wrong. We force-align each
canonical phone string to its own synthesis with an espeak-IPA CTC model,
DTW-align the hypothesis synthesis to it, and charge the spectral distortion
inside each phone's span to that phone's class. The measure is validated first:
wrong phones carry $1.72\times$ the distortion of correct ones (86.8 vs 50.5).

The verdict is unchanged. Acoustic weighting does move
in H2's favour, effect size $0.607 \rightarrow 0.756$ and ordering violations
$4/12 \rightarrow 2/12$ against the sequence metric on the same utterances,
but not far enough to clear the pre-registered threshold of $1.0$, and the
failure has the same shape as before: the ``survives'' group's internal spread
($0.120$) is $3.5\times$ the difference between the groups ($0.034$), and a
nasal again out-gains a diphthong. So ``PER mis-weights the classes'' is no
longer an available explanation for the sequence-level refutation.



One explanation does remain open, and it is not the one this test closes. Both
signals here are TTS renderings of phone strings, so the synthesiser produces
canonical formant trajectories for whatever string it is given. This measures
the cost of getting a phone's \emph{identity} wrong; it says nothing about a
wrong trajectory for a phone whose identity is right. That still needs real
converted audio against a native rendition (\S\ref{sec:remaining}). The sample
is also far smaller than the sequence-level test, 1451 phones per condition
against 295{,}785, and is weighted accordingly.

\subsection{RQ3: conversion needs more lookahead than transcription}
\label{sec:rq3}

Identical architecture, capacity, data order, seed and step count; the two arms
differ only in the target tensor. Endpoint gain $L{=}0\rightarrow640$\,ms is
$-0.2695$ (63.2\% relative) for conversion against $-0.1893$ (42.8\%) for
accent-faithful transcription, a paired difference of $+0.0802$
($\mathrm{sd}=0.0035$, $t=39.7$).

\textbf{The dense grid replicates this, and sharpens it.} Running the full
16-lookahead $\times$ 2-seed sweep on the transcription arm as well puts both
arms on the same dense grid (Table~\ref{tab:h3dense}). The endpoint gains
reproduce the 7-point result almost exactly, 63.3\% for conversion against
43.5\% for transcription, against 63.2\% and 42.8\% before, on an
independent, denser sample. The exchange rates separate cleanly:
\begin{center}
conversion $-0.0452$ vs transcription $-0.0292$ PER per doubling,
\end{center}
$R^2 = 0.983$ and $0.984$ respectively, so \textbf{conversion buys
$\mathbf{1.55\times}$ as much per doubling of lookahead}.

The mechanism is visible in the gap rather than the levels. At $L{=}0$ the two
arms are nearly the same task ($0.4256$ vs $0.4427$, a gap of $0.0171$); by
$L{=}640$\,ms the gap is $0.0939$, \textbf{5.5$\times$ wider}. It widens in
13 of 15 adjacent steps, and the transcription arm is worse at 16 of 16
lookaheads, every one by more than the $0.0044$ floor.

The arms differ in their target tensor, which carries more than the task: label
entropy, sequence length and annotation noise move with it. We therefore repeat
the test \emph{inside} the canonical-phone arm, where none of those can differ.
Splitting reference positions by whether the speaker actually deviated from
canonical ($\mathrm{g2p}\neq\mathrm{ipa}$), the exchange rate is $-0.0520$ PER
per doubling at accent-changing positions against $-0.0394$ at unchanged ones,
a factor of $\mathbf{1.32}$ from the same model on the same utterances. That
converges with the $1.55\times$ between arms. One caveat belongs with it: in
\emph{relative} terms the ordering reverses ($51\%$ against $71\%$), because
accent-changing positions start far harder ($0.590$ against $0.354$) and the
gap between the two widens with lookahead. Future context buys more absolute
accuracy where accent lives, but does not close that gap.

Two secondary points fall out. Both arms saturate at the same
${\approx}240$\,ms, and neither yields a locatable knee (breakpoints move to
$[40, 240, 160]$\,ms across axis treatments for transcription, against
$[40, 180, 180]$ for conversion), so \S\ref{sec:rq1}'s structural conclusion is
not an artefact of the canonical-phone arm. And the transcription arm's own curve is
just as log-linear, so the shallower slope is a smaller effect of lookahead, not
a noisier measurement of the same one.

\begin{table}[t]
\centering\small
\caption{Both arms on the dense grid, mean of 2 seeds. The gap is the
transcription arm's PER minus the canonical-phone arm's: near zero at $L{=}0$ and
widening monotonically with lookahead.}
\label{tab:h3dense}
\begin{tabular}{rrrr}
\toprule
$L$ (ms) & conversion & transcription & gap \\
\midrule
0   & 0.4256 & 0.4427 & $+$0.0171 \\
20  & 0.3700 & 0.3896 & $+$0.0196 \\
40  & 0.3275 & 0.3600 & $+$0.0325 \\
80  & 0.2754 & 0.3271 & $+$0.0517 \\
160 & 0.2206 & 0.2908 & $+$0.0702 \\
240 & 0.1973 & 0.2755 & $+$0.0782 \\
320 & 0.1823 & 0.2666 & $+$0.0843 \\
640 & 0.1561 & 0.2501 & $+$0.0939 \\
\midrule
\multicolumn{2}{l}{exchange rate (PER/doubling)} & & \\
\multicolumn{1}{r}{conversion} & $-0.0452$ & \multicolumn{2}{l}{$R^2=0.983$} \\
\multicolumn{1}{r}{transcription} & $-0.0292$ & \multicolumn{2}{l}{$R^2=0.984$} \\
\bottomrule
\end{tabular}
\end{table}

The claim is carried by the \emph{preference margin}, PER against the other
arm's labels minus PER against its own, which grows monotonically in 3/3
seeds for conversion ($+0.032$ at $L{=}0$ to $+0.099$ at 640\,ms; 18/18
adjacent step $\times$ seed comparisons positive) while the transcription
control is 0/3 monotone and drifts \emph{down} ($-0.0103$). Building the claim
on a difference rather than a level turned out to matter
(\S\ref{sec:lessons}).

\subsection{\tbuf{}: jitter measured, I/O not}

Two 30\,s duplex runs on the M4 at 40\,ms blocks give a required jitter buffer
of $0.14$ and $0.08$\,ms with zero over/underruns, so the jitter term is
negligible. The I/O term is \textbf{not measured}: an acoustic loopback is
defeated by platform echo cancellation, which scales with its own reference and
so cannot be overcome by a louder probe, and the driver-reported figure is
queueing rather than converter delay (its excess over the requested block is
\emph{identically} 201.69\,ms at both 10 and 20\,ms blocks). We report the gap
rather than substituting a number.

\section{A measurement failure worth reporting}
\label{sec:lessons}

\textbf{A padding bug rewrote the curve's shape.} A masked block handed a
zero-padded hidden-state tensor to a depthwise convolution and feed-forward
that do not respect the key-padding mask. The correction was not a constant
offset. It grows with lookahead, from $-0.014$\,PER at $L{=}0$ to $-0.086$
at 320\,ms, $16.4\times$ the noise floor, so the buggy curve
\emph{understated} the benefit of lookahead, biasing the central quantity of
RQ1. A constant offset would have cancelled in every within-run comparison;
this did not. Every number reported here was regenerated after the fix. Two
further failures, a breakpoint estimator that returns a breakpoint whether or
not one exists and an alphabet mismatch that silently bucketed 57\% of
reference tokens, are documented in the repository.

\section{Limitations}

\textbf{Every quality result here is phone-level, not perceptual}, and this is
a phone-translation probe rather than an end-to-end FAC system: the
latency--accuracy relation reported here need not transfer to synthesis,
prosody, speaker similarity or perceived accentedness. ``Unfunded'' would be
the wrong reason to give for the missing listening test, stimuli built from
these outputs cannot carry the judgement, because the phone sequences hold no
word or stress structure and a system at 0.16--0.42 PER emits no recognisable
words; an informal session confirmed both. Training is 1200 steps at fixed
budget, not to convergence, over a frozen encoder; the dense sweep has two
seeds per arm and no repeated cells, so its noise floor is imported from six
accidental repeats in the other sweep and carries no confidence interval of its
own. \tbuf{}'s I/O term is unmeasured, and compute is characterised on Apple
Silicon and ARM64 Linux only.

\section{Conclusion}

For causal phone translation, future context buys accuracy at an approximately
log-linear average rate of ${\approx}0.045$ PER per doubling, with no
identifiable knee and measurable returns to a saturation region of
200--280\,ms, well beyond PHONOS's 40\,ms translator budget, relative to
which a further 39.8\% reduction is available. Lookahead changes algorithmic
latency rather than per-chunk cost, though throughput must still satisfy
$\tcmp<c$ (Eq.~\ref{eq:rtf}) to run at all.

\section{Acknowledgements}

This work received no external funding and used no institutional compute: all
GPU results are from free-tier Google Colab (Tesla T4) and all CPU results from
consumer hardware. Thanks to the maintainers of L2-ARCTIC, WavLM, sherpa-onnx
and Piper, whose public releases made a single-author study of this scope
possible.

\section{Compliance with ethical standards}

This work used only the publicly released L2-ARCTIC corpus and involved no new
data collection and no human subjects, so no ethical approval was required. No
listening study was run, and no claim about perceived accentedness or
naturalness is made. The author declares no competing interests.

\bibliographystyle{IEEEbib}
\bibliography{refs}

\end{document}
